% Inclusion fragment. Requires amsmath, amssymb, amsthm and theorem,
% lemma, proposition, corollary environments. No private macros are used.
% This fragment matches Eq/D object formulas, internal Q clauses, and
% ordinary hypothesis use for object formulas only.

\section{Counter congruence recovers unrestricted equality transport}
\label{ctr:section}

The reduction in this section concerns the mathematical inference
interface. A smaller primitive interface and a smaller executable checker
are separate claims; only the accompanying measurements can establish
the latter. We use the numerical terms, the Eq/$\mathsf D$ object
formulas, the internal $\mathsf Q$ clauses, and the primitive rules
specified above. In particular,
\[
 \begin{aligned}
 R(0,a;u,v.h)&=Sa,\\
 R(Sn,a;u,v.h)&=h[n/u,R(n,a;u,v.h)/v],
 \end{aligned}
 \tag{R0, RS}
\]
where $Sa:=R(0,a;u,v.0)$.
Only the update $h$ binds $u,v$, and all substitutions are capture-free.
No variable ranges over numerical program descriptions. A template in
what follows is one fixed finite expression.

Recall the two restricted equality rules:
\[
 \frac{\Gamma\mathrel{\triangleright}s=t\qquad
       \Gamma\vdash s=r}
      {\Gamma\vdash t=r}
 \quad\mathsf{RL},
 \qquad
 \frac{\Gamma\mathrel{\triangleright}s=t}
      {\Gamma\vdash R(s,a;u,v.h)=R(t,a;u,v.h)}
 \quad\mathsf{CC}.
 \tag{Root and counter}
\]
In $\mathsf{CC}$ the seed, update, and passive parameters are the same
on both sides. The availability relation $\Gamma\mathrel{\triangleright}M$
means that $M$ is a retained assumption or an already proved object
formula. Internal $\mathsf Q$ clauses are available by membership.
Full hypothesis use is present for Eq/$\mathsf D$ object assumptions;
there is no rule with an internal $\mathsf Q$ clause as its object goal.

The proof uses $\mathsf{Hyp}$, $\mathsf{QE}$, $\mathsf{DE}$ and
$\mathsf{DI}$-case, in addition to R0, RS, $\mathsf{RL}$ and
$\mathsf{CC}$. It does not use $\mathsf{DI}$-progress. In particular,
no equation-induction instance is hidden in the reduction.

\subsection{Structural substitution is not equality transport}

Weakening, exchange, contraction, renaming, and capture-free substitution
of a numerical term for a free parameter are admissible by induction on
finite derivations. Rename eigenparameters first. R0 and RS are stable
under substitution; substitution preserves the matching seed and update
in $\mathsf{CC}$ and the matching root in $\mathsf{RL}$. It also
preserves a formula's root constructor, so a permitted progress instance
remains literal if progress occurs elsewhere in the derivation.
The ordinary eigenparameter checks handle the logical rules.
Object-assumption cut is obtained by replacing each use of the assumption
by the supplied derivation, again freshening eigenparameters.
These are operations on fixed finite proof syntax. They do not appeal
to an object-theory rule replacing equal terms inside a formula.

\subsection{Equality paths using root replacement alone}

\begin{lemma}[Equality algebra]
\label{ctr:equality-algebra}
Reflexivity, symmetry, and transitivity are derivable using RS and
$\mathsf{RL}$, without $\mathsf{CC}$ or general equality transport.
\end{lemma}
\begin{proof}
For a numerical term $s$ put
\[
             c_s:=R(Ss,0;u,v.u).
\]
RS gives $c_s=s$. Use one copy as the available major and another as
$\mathsf{RL}$'s minor; the conclusion is $s=s$.
For symmetry, an available $s=t$ and the just-derived $s=s$ give
$t=s$ by $\mathsf{RL}$.
For transitivity, suppose $s=t$ and $t=r$ are available or proved.
Hypothesis use supplies an object proof when necessary. Symmetry gives
$t=s$; apply $\mathsf{RL}$ with this major and minor $t=r$ to obtain
$s=r$. Consequently every finite path of proved numerical equations
can be composed, and any edge may be reversed. This path operation
changes no proper numerical subterm.
\end{proof}

\subsection{Numerical cases before congruence}

\begin{lemma}[Numerical cases]
\label{ctr:cases}
For any fixed Eq/$\mathsf D$ object template $A(y)$, the following rule
is admissible without either equality rule or progress introduction:
\[
 \frac{\Gamma\vdash A(0)\qquad\Gamma\vdash A(Sz)}
      {\Gamma\vdash A(t)}.
 \tag{Cases}
\]
Here $z$ is fresh for the retained context, the fixed parameters of
$A$, and the target term $t$; it may always be renamed to meet these
conditions.
\end{lemma}
\begin{proof}
Choose a bound coordinate $x$ absent from $A$. By weakening, the
successor premise remains available under the extra assumption
$A(0)$. Apply $\mathsf{DI}$-case with witness zero to obtain
$\Gamma\vdash\mathsf D_{x,y}A(y)$. Eliminate this paired formula.
Because $x$ is unused, its new internal clause is exactly
$\mathsf Q_yA(y)$; the fresh existential parameter does not occur in
that clause. Weaken the given base proof into the extended context and
apply $\mathsf{QE}$ at $t$. This proves $A(t)$ there, and
$\mathsf{DE}$ discharges the clause. There was no rewrite by an
equality and no induction hypothesis.
\end{proof}

\subsection{Representing a numerical context through the counter}

Define
\[
                  W(n):=R(n,0;u,v.u).
\]
The computations $W(0)=S0$ and $W(Sn)=n$ are instances of R0 and RS.
In particular the latter identity is available before any general
congruence principle has been established.

\begin{lemma}[Counter representation of a fixed context]
\label{ctr:context-representation}
For any fixed numerical template $F(x)$, choose all new bound variables
fresh for its passive parameters and put
\[
 \begin{aligned}
  H_F(n)&:=R(n,F(0);u,v.S(F(Su))),\\
  J_F(n)&:=W(H_F(n)).
 \end{aligned}
 \tag{Counter representation}
\]
Then the restricted calculus proves
\[
                         J_F(n)=F(n)
 \tag{Representation}\label{ctr:representation-equation}
\]
for arbitrary numerical $n$ and arbitrary retained context. The proof
uses numerical cases and no progress introduction.
\end{lemma}
\begin{proof}
The update of $H_F$ does not contain its state variable $v$.
Consequently R0 and RS give the following identities literally, with
no equality substitution into an enclosing expression:
\[
 H_F(0)=S(F(0)),\qquad H_F(Sz)=S(F(Sz)).
 \tag{Two computations}
\]
For either of these instances, write its argument as $q$, where $q=0$
or $q=Sz$. Apply $\mathsf{CC}$ to the displayed equality using seed
zero and the update that returns its old counter. This produces
\[
                 W(H_F(q))=W(S(F(q))).
\]
The right-hand expression has the computation equation
$W(S(F(q)))=F(q)$. Compose the two equation paths by
Lemma~\ref{ctr:equality-algebra}; the result is
$J_F(q)=F(q)$. We have now supplied both premises of
Lemma~\ref{ctr:cases} for the literal equation
$J_F(y)=F(y)$. Cases proves
\eqref{ctr:representation-equation} at an arbitrary $n$.

Every proper-subterm change in this argument is the single, explicitly
stated $\mathsf{CC}$ instance through $W$. The remaining equation
compositions use root replacement. In particular one must not silently
replace $H_F(q)$ by $S(F(q))$ under $W$ before applying
$\mathsf{CC}$; that is precisely the congruence being justified.
\end{proof}

The displayed $J_F$ is a specialization of the successor-recursion
compiler already present in the inspected drafts: take its initializer
to be $F(0)$ and its step to be $h(u,z)=F(Su)$, which ignores the state.
The additional observation here is that this specialization is proved
correct using only root replacement and congruence in the counter.
The general recursion compiler itself is inherited.

\begin{theorem}[All numerical congruence from counter congruence]
\label{ctr:all-congruence}
For every fixed numerical template $F(x)$, the restricted calculus admits
\[
 \frac{\Gamma\mathrel{\triangleright}s=t}
      {\Gamma\vdash F(s)=F(t)}.
 \tag{Numerical congruence}
\]
\end{theorem}
\begin{proof}
Use the major $s=t$ in $\mathsf{CC}$ with seed $F(0)$ and update
$S(F(Su))$. Its conclusion is $H_F(s)=H_F(t)$. A second
$\mathsf{CC}$ instance, now for $W$, gives
$J_F(s)=J_F(t)$. Lemma~\ref{ctr:context-representation} supplies
$J_F(s)=F(s)$ and $J_F(t)=F(t)$. Symmetry and path composition from
Lemma~\ref{ctr:equality-algebra} give
\[
                F(s)=J_F(s)=J_F(t)=F(t).
\]
These are equations joined at their roots. Neither the proof nor the
statement assumes an interpretation limited to standard numerals.
\end{proof}

\begin{corollary}[Arbitrary atomic transport]
\label{ctr:atomic-transport}
For all fixed numerical templates $U(x),V(x)$,
\[
 \frac{\Gamma\mathrel{\triangleright}s=t\qquad
       \Gamma\vdash U(s)=V(s)}
      {\Gamma\vdash U(t)=V(t)}
 \tag{Atomic transport}
\]
is admissible.
\end{corollary}
\begin{proof}
Theorem~\ref{ctr:all-congruence} gives $U(s)=U(t)$ and $V(s)=V(t)$.
Reverse the first equality and compose only at equation roots:
\[
                U(t)=U(s)=V(s)=V(t).
\]
\end{proof}

\subsection{A constant bound on inference-node overhead}

Count one node for each primitive inference, including a hypothesis
lookup, and count a proof tree with repeated subproofs repeated. In the
explicit expansion supplied with the report, reflexivity has three
nodes. Symmetry adds four nodes to its supplied proof, and transitivity
adds five nodes to the sum of its two supplied proofs.

Each of the zero and successor proofs of $J_F(q)=F(q)$ has eight
nodes: a computation, one CC lifting, a second computation, and the
five-node transitivity overhead. The cases construction consists of
DE, DI-case, QE, one clause hypothesis lookup, and these two eight-node
proofs. Thus $J_F(t)=F(t)$ has twenty nodes.

If the supplied proof of $s=t$ has $P$ nodes, the two CC liftings have
$P+2$ nodes, the reversed left representation has $20+4$ nodes,
and the right representation has twenty. Two transitivity constructions
therefore give
\[
                  (P+2)+(20+4)+20+5+5=P+56.
\]
The additional inference-node count is independent of the syntactic
depth of the fixed numerical template $F$. Freshening binders changes
names, not the number of inferences. This is not a constant byte-size
or running-time bound: the proof's term and formula labels contain
copies of $F$ and of the supplied endpoints. It is also not a constant
bound for the subsequent compound-formula lifting. The delivered
successor, initializer, and update-body congruence certificates each
have 57 nodes when their supplied equality is a one-node computation.

\subsection{Compound formulas and internal clauses}

The following lifting pattern is inherited from the earlier
fixed-position and atomic-transport arguments. It is included to show
that the numerical result recovers the full logical interface, and to
make its independence from compound induction explicit.

\begin{theorem}[Full object-formula transport]
\label{ctr:full-transport}\label{thm:atomic-transport}
For every fixed Eq/$\mathsf D$ object template $A(w)$, the rule
\[
 \frac{\Gamma\mathrel{\triangleright}s=t\qquad
       \Gamma\vdash A(s)}
      {\Gamma\vdash A(t)}
 \tag{Object transport}
\]
is admissible. Its expansion introduces no progress instances.
\end{theorem}
\begin{proof}
Induct externally on the number of paired binders in $A$. The induction
statement is uniform in all retained contexts, source terms, and passive
parameters. For a literal equation, use
Corollary~\ref{ctr:atomic-transport}.

Suppose $A(w)=\mathsf D_{x,y}B(w,x,y)$. Rename the two bound
variables away from $s,t$ and the retained context. Eliminate the given
$A(s)$ at a fresh parameter $a$. The extended context retains
\[
                     \mathsf Q_y B(s,a,y).
\]
To obtain the target $\mathsf D_{x,y}B(t,x,y)$, use
$\mathsf{DI}$-case with witness $a$ and a fresh successor parameter
$z$. Its local object assumption is $B(t,a,0)$; hypothesis use gives
an object proof of it. Symmetry gives $t=s$. The external induction
hypothesis applied backwards to the smaller template $B(w,a,0)$
therefore gives $B(s,a,0)$. Specialize the retained clause at $Sz$ to
obtain $B(s,a,Sz)$. Apply the induction hypothesis forwards along
$s=t$, now to $B(w,a,Sz)$, obtaining $B(t,a,Sz)$. This is exactly
the required case-introduction premise. Close $\mathsf{DI}$-case and
then $\mathsf{DE}$.

The new parameter $z$ is absent from the retained clause, $a$, and all
fixed matrix parameters; the eliminated witness $a$ is absent from the
outer conclusion. Both recursive transport uses concern the strictly
smaller object matrix $B$. This is recursion on finite external syntax,
not an induction rule of the arithmetic. The only introductions added
in the binder case are case introductions.
\end{proof}

\begin{proposition}[Transporting the use of an internal clause]
\label{ctr:clause-transport}
An available clause $\mathsf Q_yB(s,y)$, together with an available
$s=t$, can be used with exactly the operational force of
$\mathsf Q_yB(t,y)$: from any object proof of $B(t,0)$ one can
obtain $B(t,r)$ for any numerical $r$.
\end{proposition}
\begin{proof}
Transport the supplied base backwards along $t=s$, specialize the old
clause at $r$, and transport the result forwards along $s=t$, using
Theorem~\ref{ctr:full-transport} for the two object matrices.
This statement does not add a rule proving an internal clause as an
object formula. If an object representative is required, use
$\mathsf D_{x,y}B(t,y)$ with unused $x$; case introduction and the
just-described clause use construct that representative.
\end{proof}

\begin{corollary}[Exact identity-language equivalence]
\label{ctr:exact-equivalence}
Let $\mathcal T_{\mathrm{ctr}}$ have R0, RS, object hypothesis use,
$\mathsf{RL}$, $\mathsf{CC}$, $\mathsf{QE}$,
$\mathsf{DE}$, and the case and literal-equation progress introductions.
Let $\mathcal T_{\mathrm{tr}}$ instead use unrestricted Eq/$\mathsf D$
object transport as its equality rule, with all other numerical and
logical syntax and rules the same. For every finite context of object
assumptions and internal clauses, and every object formula $A$,
\[
 \Gamma\vdash_{\mathcal T_{\mathrm{ctr}}}A
 \quad\Longleftrightarrow\quad
 \Gamma\vdash_{\mathcal T_{\mathrm{tr}}}A.
\]
The translations are effective on finite derivations and preserve their
end sequents literally.
\end{corollary}
\begin{proof}
From right to left, translate premise derivations recursively and replace
each unrestricted transport instance by
Theorem~\ref{ctr:full-transport}. All other rules are copied.
From left to right, $\mathsf{RL}$ is the unrestricted-transport instance
with motive $w=r$. Reflexivity is derived from a computation equation
by that same root instance. For $\mathsf{CC}$, begin with reflexivity
$R(s,a;h)=R(s,a;h)$ and transport its left occurrence of $s$ to $t$;
this gives $R(t,a;h)=R(s,a;h)$. Derived symmetry gives the required
orientation. The other rules are again copied. All side conditions are
preserved by freshening local parameters.
\end{proof}

Because this equivalence leaves the native language and all primitive
meanings unchanged, the whole-domain bi-interpretations and actual-equality
composite isomorphisms of $\mathcal T_{\mathrm{tr}}$ transfer to
$\mathcal T_{\mathrm{ctr}}$. A proof of that arithmetic and logical
bridge is supplied separately; exact sequent equivalence is not a
substitute for stating the bridge. The present argument supplies no
absolute lower bound on checker size and no claim of novelty outside
the specifically inspected sources.
